\begin{problem}[第34届物理竞赛复赛][电磁感应]{75}
    如俯视图，在水平面内有两个分别以 $O$ 点与 $O_{1}$ 点为圆心的导电半圆弧内切于 $M$ 点，半圆 $O$ 的半径为 $2a$，半圆 $O_{1}$ 的半径为 $a$；两个半圆弧和圆 $O$ 的半径 $ON$ 围成的区域内充满垂直于水平面向下的匀强磁场（未画出），磁感应强度大小为 $B$；其余区域没有磁场。半径 $OP$ 为一均匀细金属棒，以恒定的角速度 $\omega$ 绕 $O$ 点顺时针旋转，旋转过程中金属棒 $OP$ 与两个半圆弧均接触良好。已知金属棒 $OP$ 的电阻为 $R$，两个半圆弧的电阻可忽略。开始时 $P$ 点与 $M$ 点重合。在 $t$（$0 \leq t \leq \frac{\pi}{\omega}$）时刻，半径 $OP$ 与半圆 $O_{1}$ 交于 $Q$ 点。求
    \begin{subproblem}
        沿回路 $QPMQ$ 的感应电动势；
    \end{subproblem}
    \begin{subproblem}
        金属棒 $OP$ 所受到的原磁场 $B$ 的作用力的大小。
    \end{subproblem}
\end{problem}